\title{DeepSeekMath: Pushing the Limits of Mathematical Reasoning in Open Language Models}

\begin{abstract}
Mathematical reasoning remains a formidable task for language models because of its intricate and highly structured characteristics. This work presents DeepSeekMath~7B, a model created by further pre-training DeepSeek-Coder-Base-v1.5 7B on a dataset comprising 120B math-centric tokens collected from Common Crawl, in combination with text and programming code. DeepSeekMath~7B attains a notable 51.7\% accuracy on the MATH benchmark at competition level, achieving results without the need for external tools or voting mechanisms and closely matching the proficiency of Gemini-Ultra and GPT-4. When utilizing self-consistency with 64 generated samples, DeepSeekMath~7B reaches a 60.9\% score on MATH. The heightened mathematical reasoning capacity of DeepSeekMath~is attributed to two principal components: First, an advanced data selection process systematically taps into the extensive publicly accessible web data. Second, we develop Group Relative Policy Optimization (GRPO), an adaptation of Proximal Policy Optimization (PPO), which not only sharpens mathematical reasoning but also improves PPO’s memory efficiency.
\end{abstract}

\section{Introduction}

The development of large language models (LLMs) has fundamentally transformed methodologies in mathematical reasoning within AI, yielding remarkable progress across quantitative reasoning assessments~\citep{MATH} and geometry benchmarks~\citep{trinh2024solving}. Additionally, these systems have demonstrated their utility in supporting human efforts to solve intricate mathematical tasks~\citep{tao}. Nevertheless, state-of-the-art systems like GPT-4~\citep{gpt4} and Gemini-Ultra~\citep{gemini} remain proprietary, and the existing suite of open-source models still lags substantially in terms of capability and performance.

This paper presents DeepSeekMath, a specialized language model that demonstrates markedly enhanced mathematical reasoning compared to existing open-source models and attains near-GPT-4 performance on a range of academic tasks. The advancement is made possible by the construction of the DeepSeekMath Corpus, an extensive and high-quality pre-training dataset containing 120B math tokens. The corpus is sourced from Common Crawl (CC) through the utilization of a classifier based on fastText \citep{joulin2016fasttext}. In its first stage, the classifier is trained with positive samples from OpenWebMath \citep{openwebmath} and a heterogeneous set of other web content as negative examples. Following this, the trained classifier is leveraged to extract further positive examples from the CC, which are subsequently subjected to refinement via human annotation. The classifier is retrained with this improved dataset, resulting in performance gains. Empirical results reveal the high quality of the corpus, with DeepSeekMath-Base 7B attaining 64.2\% accuracy on GSM8K \citep{gsm8k} and 36.2\% on the competition-level MATH benchmark \citep{MATH}, surpassing the performance of Minerva 540B \citep{minerva}. Furthermore, owing to the multilingual design of the DeepSeekMath Corpus, we observe gains on Chinese mathematics benchmarks \citep{wei2023cmath,agieval}. We propose that the methodologies developed for mathematical data curation in this project can serve as a valuable foundation for subsequent research, with considerable potential for future refinement and advancement.

DeepSeekMath-Base is constructed by initializing from DeepSeek-Coder-Base-v1.5 7B \citep{deepseek-coder}, as we find that leveraging a model pre-trained on code outperforms beginning with a generic LLM. Additionally, math-oriented training confers improvements not only in mathematical tasks but also yields gains on broader reasoning benchmarks, such as MMLU \citep{mmlu} and BBH \citep{bbh}, suggesting a transfer of enhanced reasoning capabilities.

Post pre-training, we refine DeepSeekMath-Base through mathematical instruction tuning using chain-of-thought \citep{cot}, program-of-thought \citep{pot,pal}, and reasoning integrated with tools \citep{tora}. The resultant DeepSeekMath-Instruct 7B surpasses all existing 7B models and demonstrates performance on par with leading 70B open-source models that have undergone instruction tuning.

In addition, we propose Group Relative Policy Optimization (GRPO), a novel RL algorithm adapted from Proximal Policy Optimization (PPO) \citep{schulman2017proximal}. Unlike PPO, GRPO omits the need for a critic model, instead utilizing group-based score baselines, which leads to notable reductions in computational requirements. By employing only a subset of English instruction tuning data, GRPO delivers substantial performance increases over the already robust DeepSeekMath-Instruct, as shown by marked improvements in both within-domain (GSM8K: 82.9\% $\rightarrow$ 88.2\%, MATH: 46.8\% $\rightarrow$ 51.7\%) and out-of-domain mathematical datasets (e.g., CMATH: 84.6\% $\rightarrow$ 88.8\%) in the RL phase. We further introduce a unifying framework to contextualize multiple approaches, including Rejection Sampling Fine-Tuning (RFT) \citep{yuan2023scaling}, Direct Preference Optimization (DPO) \citep{dpo}, PPO, and GRPO. Within this unified framework, we observe that these approaches fundamentally align as forms of direct or simplified RL. Comprehensive ablation studies are conducted, such as comparisons of online vs. offline training, outcome vs. process supervision, and single-turn vs. iterative RL, to systematically explore the core properties of this framework. Finally, we elucidate how RL drives advancements in instruction-tuned models and outline possible pathways for realizing more efficient RL under this unified perspective.

\subsection{Contributions}
This work presents advancements in large-scale math-focused pre-training as well as an in-depth examination of reinforcement learning methods.

\textbf{Math Pre-Training at Scale}

\begin{itemize}[topsep=0pt]
    \item We demonstrate that Common Crawl, an open dataset, offers a wealth of mathematical material. By developing a systematic pipeline for data selection, we curate the DeepSeekMath Corpus—comprising 120B high-quality, math-relevant tokens extracted from web sources. This corpus is nearly seven times larger than the dataset employed by Minerva \citep{minerva} and nine times greater than OpenWebMath's \citep{openwebmath} collection.
    
    \item Our DeepSeekMath-Base 7B pre-trained model matches the performance level of the substantially larger Minerva 540B model \citep{minerva}, suggesting that sheer model scale is not the sole determinant for mathematical reasoning proficiency. These results highlight that a comparatively modest-sized model, when trained on well-curated data, can achieve exceptional outcomes.
    
    \item We present empirical observations from our math pre-training experiments, noting that exposure to code data before mathematical tasks leads to enhanced problem-solving performance in mathematical contexts, both with and without the use of auxiliary tools. This partially addresses the persistent query: \textit{does code-related pre-training boost reasoning skills?} Our results indicate an affirmative, at least in mathematical domains.
    
    \item While using arXiv papers as a training source is standard in math-related machine learning research, we find no significant benefit across the mathematical benchmarks evaluated herein.
\end{itemize}

\textbf{Exploration and Analysis of Reinforcement Learning}

\begin{itemize}[topsep=0pt]
    \item We present Group Relative Policy Optimization (GRPO), a reinforcement learning algorithm characterized by both high efficiency and effectiveness. GRPO eliminates the need for a critic network by estimating baselines directly from group scores, thereby substantially lowering the computational demands required for training in comparison to Proximal Policy Optimization (PPO).
    \item Our results show that GRPO notably improves the effectiveness of the instruction-tuned model DeepSeekMath-Instruct when relying solely on instruction-tuning datasets. In addition, we find that reinforcement learning with GRPO leads to better generalization on out-of-domain tasks.
    \item We introduce a cohesive framework that unifies the understanding of various approaches, including RFT, DPO, PPO, and GRPO. Through extensive empirical studies—including comparisons of online and offline learning, supervision based on outcomes versus processes, and contrasts between single-turn and iterative reinforcement learning—we thoroughly analyze the critical aspects underpinning this framework.
    \item Drawing on insights from this unified perspective, we examine the mechanisms that underlie the success of reinforcement learning and highlight several promising research directions for advancing large language model (LLM) reinforcement learning.
\end{itemize}

\subsection{Summary of Evaluations and Metrics}

\begin{itemize}[topsep=0pt]
    \item \textbf{English and Chinese Mathematical Reasoning}: Our evaluation includes a thorough analysis of model performance across both English and Chinese datasets, encompassing mathematical questions ranging from elementary to collegiate difficulty. For English, the benchmarks comprise GSM8K \citep{gsm8k}, MATH \citep{MATH}, SAT \citep{llemma}, OCW Courses \citep{minerva}, and MMLU-STEM \citep{mmlu}. Chinese evaluations draw on MGSM-zh \citep{mgsm}, CMATH \citep{wei2023cmath}, Gaokao-MathCloze \citep{agieval}, and Gaokao-MathQA \citep{agieval}. We assess the ability of our models to produce comprehensive textual solutions without auxiliary tools, as well as their problem-solving proficiency using Python.

    On English tasks, DeepSeekMath-Base demonstrates performance on par with the proprietary Minerva 540B \citep{minerva}, while significantly outperforming open-source base models such as Mistral 7B \citep{mistral} and Llemma-34B \citep{llemma}—regardless of their exposure to mathematics-specific pre-training—frequently with substantial leads. Remarkably, DeepSeekMath-Base also leads on Chinese tasks. This may be attributed to our strategy of incorporating high-quality, diverse-language data, in contrast to the English-centric data curation in previous studies \citep{minerva,llemma}. Further, following instruction tuning and reinforcement learning, both DeepSeekMath-Instruct and DeepSeekMath-RL yield impressive results, for the first time exceeding 50\% accuracy on the challenging, competition-grade MATH benchmark within the open-source domain.
\end{itemize}

\item \textbf{Formal Mathematics}: For assessing DeepSeekMath-Base, we utilize the informal-to-formal theorem proving task from \citep{dsp_proof} applied to the miniF2F benchmark \citep{minif2f}, with Isabelle \citep{isabelle} designated as the proof assistant. DeepSeekMath-Base achieves notable few-shot autoformalization results in this setting.

\item \textbf{Natural Language Understanding, Reasoning, and Code}: To capture a broad spectrum of the model’s proficiency in general comprehension, logical reasoning, and programming, we evaluate DeepSeekMath-Base on the Massive Multitask Language Understanding (MMLU) benchmark \citep{mmlu}, which consists of 57 multiple-choice problems spanning a wide array of topics. Additionally, we use BIG-Bench Hard (BBH) \citep{bbh}, which features 23 difficult tasks requiring predominantly multi-step reasoning, as well as HumanEval \citep{codex} and MBPP \citep{mbpp}, both standard benchmarks for measuring code generation abilities. The pre-training process on mathematical data enhances both the language understanding and reasoning capacities of the model.

\section{Math Pre-Training}

\subsection{Data Collection and Decontamination} 
This section details the methodology for assembling the DeepSeekMath Corpus utilizing data from Common Crawl. Figure \ref{fig:data_collect} illustrates our iterative approach, which systematically scales up a mathematical text corpus sourced from Common Crawl, beginning with a high-quality but relatively limited seed dataset focused on mathematics. Notably, this strategy extends beyond mathematics and is applicable to areas like programming as well.

The initial seed set is OpenWebMath \citep{openwebmath}, recognized for its curated mathematical web resources. Using this foundation, we train a fastText model \citep{joulin2016fasttext} to identify additional web pages with similar mathematical content. To construct training data, 500,000 entries from the seed are randomly sampled as positive examples, while an equal number of Common Crawl web pages serve as negatives. Model training utilizes an open-source implementation\footnote{\url{https://fasttext.cc}}, with parameters set as follows: a vector size of 256, learning rate of 0.1, maximum n-gram length of 3, minimum word frequency of 3, and three training epochs. To further streamline Common Crawl, we employ both exact (URL-based) and approximate (near-duplicate) deduplication, narrowing it down to 40B HTML pages. The fastText model is then used to extract mathematics-relevant pages from this deduplicated pool. To enhance quality, all candidates are scored according to fastText’s predictions, and only those ranked highest are selected for further processing. Subsequent pre-training experiments are conducted on datasets containing the top 40B, 80B, 120B, and 160B tokens, with the first round focusing on retaining the upper 40B tokens.

Following the initial phase of data gathering, a significant number of mathematical web pages remain uncollected. This issue primarily arises from the lack of diversity among positive examples used to train the fastText model. To address this, we enhance the seed dataset by identifying and incorporating additional mathematical web sources, thereby refining the fastText model. Our method begins with partitioning the entire Common Crawl dataset by domain, where a domain refers to all web pages sharing the same root URL. For each domain, we compute the proportion of web pages captured during the initial collection. Domains with more than 10\% of their pages collected are designated as math-related (for instance, \url{mathoverflow.net}).

Next, we manually review and annotate the URLs within these domains that pertain to mathematical topics (such as \url{mathoverflow.net/questions}). Pages linked to these annotated URLs, which were previously missed, are incorporated into the seed dataset. By enriching the set of positive examples in this way, we facilitate the development of a more robust fastText model that is better equipped to retrieve mathematical data in future collection cycles. After completing four cycles of data gathering, we acquire 35.5M mathematical web pages, encompassing a total of 120B tokens. By the fourth round, we observe that approximately 98\% of the pages had already been gathered during the previous iteration, prompting us to halt further data collection.

To prevent benchmark contamination, we employ the filtering protocol detailed in \cite{deepseek-coder}. Specifically, web pages that contain questions or answers originating from English mathematical benchmarks, such as GSM8K \citep{gsm8k} and MATH \citep{MATH}, or Chinese benchmarks, including CMATH \citep{wei2023cmath} and AGIEval \citep{agieval}, are systematically removed. The filtering procedure involves eliminating any text fragment containing a 10-gram substring that exactly matches any portion of the evaluation benchmarks. For benchmark excerpts shorter than 10 grams but at least 3 grams in length, we conduct exact string matching to exclude contaminated entries from the mathematical training set.

\subsection{Evaluating the DeepSeekMath~Corpus Quality}

To assess the quality of the DeepSeekMath~Corpus, we conduct pre-training experiments comparing it with other recent large-scale mathematical datasets:

\begin{itemize}[topsep=0pt]
    \item \textbf{MathPile} \citep{mathpile}: This multi-source dataset (8.9B tokens) is compiled from materials such as textbooks, Wikipedia, ProofWiki, CommonCrawl, StackExchange, and arXiv, with more than 85\% of the content sourced from arXiv;
    \item \textbf{OpenWebMath} \citep{openwebmath}: A 13.6B token dataset derived from CommonCrawl through targeted filtering for mathematical material;
    \item \textbf{Proof-Pile-2} \citep{llemma}: A composite corpus that brings together OpenWebMath, AlgebraicStack (10.3B tokens of math-related code), and arXiv papers (28.0B tokens). For Proof-Pile-2 experiments, we adhere to the arXiv:Web:Code ratio of 2:4:1 as described in \cite{llemma}.
\end{itemize}

\subsubsection{Training Setting}
\label{sec:quality-policy}
We fine-tune a general-purpose pre-trained language model consisting of 1.3B parameters—utilizing the architecture identical to the DeepSeek LLMs \citep{deepseek-llm}, hereafter referred to as DeepSeek-LLM 1.3B—on distinct mathematical corpora, allocating 150B tokens for each dataset. Training across all models is implemented via the streamlined and resource-efficient HAI-LLM \citep{haillm} framework. Adhering to the protocols established by DeepSeek LLMs, we employ the AdamW optimizer \citep{adamW} with hyperparameters set to $\beta_1=0.9$, $\beta_2=0.95$, and a $\mathrm{weight\_decay}$ of 0.1. The learning rate follows a multi-phase schedule: initially warming up over 2,000 steps to the maximum rate, subsequently dropping to 31.6\% at the 80\% mark of total steps, and finally tapering to 10.0\% after 90\% of training has concluded. The peak learning rate is configured at 5.3e-4, and we utilize a batch size equivalent to 4 million tokens, each sample comprising a 4K token context window.

\subsubsection{Evaluation Results}

\textbf{The DeepSeekMath~Corpus stands out for its superior quality, extensive multilingual content, and largest available scale.}

\begin{itemize}[topsep=0pt]
    \item \textbf{High-quality}:
    We assess downstream task performance across eight mathematical benchmarks using few-shot chain-of-thought prompting \cite{cot}. Results in Table \ref{tab:corpora_comparison} indicate that models trained with DeepSeekMath~Corpus consistently outperform alternatives. Furthermore, Figure \ref{fig:corpora_comparisons} illustrates that at 50B tokens (one full epoch over Proof-Pile-2), the DeepSeekMath-trained model already surpasses that trained on Proof-Pile-2, highlighting the corpus’s higher average sample quality.
    \item \textbf{Multilingual}:
    The DeepSeekMath~Corpus includes data from numerous languages, with English and Chinese comprising the largest fractions. Table \ref{tab:corpora_comparison} shows that leveraging the DeepSeekMath~Corpus leads to notable gains in mathematical reasoning tasks for both English and Chinese. Conversely, models trained on existing corpora—which are mostly centered on English—exhibit only marginal, if any, advances for Chinese, and may even hinder mathematical reasoning performance in this language.
    \item \textbf{Large-scale}:
    The DeepSeekMath~Corpus substantially exceeds the size of previously available mathematical datasets. As seen in Figure \ref{fig:corpora_comparisons}, DeepSeek-LLM 1.3B models trained on this larger dataset exhibit a significantly steeper learning curve and more sustained performance improvements over time. In comparison, baseline datasets are considerably smaller and are cycled through multiple times during training, causing their learning progress to plateau early.
\end{itemize}

\subsection{Training and Evaluating DeepSeekMath-Base 7B}

This section presents DeepSeekMath-Base 7B, a foundational model that demonstrates strong mathematical reasoning proficiency. The initial model weights are derived from DeepSeek-Coder-Base-v1.5 7B \citep{deepseek-coder}, and subsequent training is carried out over 500B tokens. The composition of the training dataset is as follows: 56\% from the DeepSeekMath Corpus, 4\% sourced from AlgebraicStack, 10\% from arXiv publications, 20\% consists of Github code repositories, while the final 10\% contains natural language data extracted from Common Crawl in both English and Chinese. The majority of training parameters align with those outlined in Section \ref{sec:quality-policy}, but with two modifications: the peak learning rate is adjusted to 4.2e-4, and the token batch size is increased to 10M.

To thoroughly gauge DeepSeekMath-Base 7B's mathematical prowess, we examine the model's capability to independently construct mathematical solutions, utilize computational tools for problem solving, and undertake formalized theorem proving. We further supplement these evaluations with a broader analysis of the base model's general competencies, specifically natural language understanding, deductive reasoning, and programming ability.

\paragraph{Stepwise Mathematical Problem Solving}

The evaluation framework involves measuring DeepSeekMath-Base's skill in mathematical problem solving using few-shot chain-of-thought prompting \citep{cot}, spanning eight different benchmarks in both English and Chinese. These include datasets targeting quantitative reasoning, such as GSM8K \citep{gsm8k}, MATH \citep{MATH}, and CMATH \citep{wei2023cmath}, as well as multiple-choice tasks like MMLU-STEM \citep{mmlu} and Gaokao-MathQA \citep{agieval}. Collectively, these benchmarks represent a diverse range of mathematical domains and levels, extending from primary to undergraduate curricula.

As detailed in Table \ref{tab:base_math_cot}, among publicly available base models, DeepSeekMath-Base 7B achieves the highest scores on all eight benchmarks, outperforming widely adopted models such as Mistral 7B \citep{mistral} and the more recent Llemma 34B \citep{llemma}, the latter of which received specialized math training via Proof-Pile-2 \citep{llemma}. Remarkably, on the highly challenging MATH benchmark, DeepSeekMath-Base delivers an absolute improvement of over 10\% relative to existing open-source base models, and even exceeds the performance of Minerva 540B \citep{minerva}—a proprietary model that is 77 times larger, based on PaLM \citep{palm}, and extensively trained on mathematical corpora.

\paragraph{Mathematical Problem Solving with Tool Use} We assess the ability of models to perform program-assisted mathematical reasoning using the GSM8K and MATH datasets, employing few-shot program-of-thought prompting strategies \citep{pot,pal}. Each model receives prompts instructing them to compose Python programs (utilizing libraries such as \textit{math} and \textit{sympy} for complex calculations) as solutions, and their outputs are considered for final evaluation. As presented in Table \ref{tab:base_math_pot_proof}, DeepSeekMath-Base 7B surpasses the previously leading open-source model, Llemma 34B, on these program-based tasks.

\paragraph{Formal Mathematics}

Automating the generation of formal mathematical proofs both promotes accuracy and bolsters the reliability of proofs, leading to growing interest in this research direction. We benchmark DeepSeekMath-Base 7B on the informal-to-formal proof conversion task introduced in \citep{dsp_proof}, which requires producing a formal proof given an informal problem statement, its formal equivalent, and an informal proof outline. Evaluations are conducted on miniF2F \citep{minif2f}, a dataset centered on Olympiad-level mathematics, by tasking the model to output formal proofs in Isabelle under few-shot prompt settings. Following the methodology of \cite{dsp_proof}, generated proof sketches are post-processed with Sledgehammer \citep{sledgehammer}, an automated proving assistant, to complete any omitted proof steps. Results in Table \ref{tab:base_math_pot_proof} indicate that DeepSeekMath-Base 7B achieves strong results in proof formalization.

\paragraph{Natural Language Understanding, Reasoning, and Code}

We measure performance in natural language understanding with MMLU \citep{mmlu}, reasoning skills with BBH \citep{bbh}, and programming on HumanEval \citep{codex} and MBPP \citep{mbpp}. Table \ref{tab:understanding_reasoning_code} shows that DeepSeekMath-Base 7B markedly improves its results on both MMLU and BBH compared to its predecessor, DeepSeek-Coder-Base-v1.5 \citep{deepseek-coder}, highlighting the efficacy of mathematical training for language and reasoning abilities. Furthermore, the inclusion of code token training enables DeepSeekMath-Base 7B to preserve the code generation performance observed in DeepSeek-Coder-Base-v1.5 on HumanEval and MBPP. Collectively, DeepSeekMath-Base 7B demonstrates substantial superiority over the general-purpose Mistral 7B \citep{mistral} across the evaluated reasoning and code benchmarks.

\section{Supervised Fine-Tuning}

\subsection{SFT Data Curation}

We assemble a dataset for mathematical instruction tuning that encompasses both English and Chinese mathematical problems, sourced from a range of domains and spanning a spectrum of difficulty levels. Each problem is aligned with its solution in one of the following structured formats: chain-of-thought (CoT) \citep{cot}, program-of-thought (PoT) \citep{pot,pal}, or through the use of integrated computational tools as per \citep{tora}. The training corpus consists of 776,000 examples.

\begin{itemize}[topsep=0pt]
    \item \textbf{English mathematical datasets}: We enhance datasets such as GSM8K and MATH by providing tool-integrated solution annotations, incorporate a filtered selection from MathInstruct \citep{MathInstruct}, and include Lila-OOD’s training split \citep{lila} where solutions adopt either CoT or PoT methodologies. The English component of our dataset spans multiple branches of mathematics, including but not limited to algebra, probability, number theory, calculus, and geometry.
    \item \textbf{Chinese mathematical datasets}: Our Chinese-language data consists of K-12 level math questions, covering 76 distinct subcategories (e.g., linear equations), each accompanied by CoT and tool-augmented reasoning solutions.
\end{itemize}

\subsection{Training and Evaluating DeepSeekMath-Instruct 7B}

Here, we detail the supervised fine-tuning process applied to DeepSeekMath-Instruct 7B, which builds upon DeepSeekMath-Base via specialized mathematical instruction tuning. For each training batch, input samples are concatenated at random until a maximum sequence length of 4,000 tokens is reached. The model is then trained over 500 iterations, using a batch size of 256 and maintaining a fixed learning rate of $5 \times 10^{-5}$.

The assessment of our models’ mathematical capabilities, both in tool-free and tool-assisted modes, is performed on four established quantitative reasoning benchmarks available in English and Chinese. We compare the results with those from top-performing contemporary systems:

\begin{itemize}[topsep=0pt]
    \item \textbf{Closed-source models}: These comprise (1) the GPT series, notably GPT-4 \citep{gpt4} and GPT-4 Code Interpreter~\footnote{\url{https://openai.com/blog/chatgpt-plugins\#\#code-interpreter}}; (2) Gemini Ultra and Pro \citep{gemini}; (3) Inflection-2 \citep{inflection-2}; (4) Grok-1~\footnote{\url{https://x.ai/model-card}}; alongside recently unveiled models from Chinese developers including (5) Baichuan-3~\footnote{\url{https://www.baichuan-ai.com}} and (6) GLM-4~\footnote{\url{https://open.bigmodel.cn/dev/api\#glm-4}} from the GLM lineage \citep{glm}.
\end{itemize}

These models are intended for broad applications, with the majority having been subjected to multiple alignment processes.

\item \textbf{Open-source models} featured in this work include: general-purpose models such as (1) DeepSeek-LLM-Chat 67B \citep{deepseek-llm}, (2) Qwen 72B \citep{qwen}, (3) SeaLLM-v2 7B \citep{seallm}, and (4) ChatGLM3 6B \citep{chatglm3}. We also consider models with specific enhancements in mathematical capability, such as (5) InternLM2-Math 20B~\footnote{\url{https://github.com/InternLM/InternLM-Math}}, an InternLM2-based model refined through math-targeted training and subsequent instruction tuning, (6) Math-Shepherd-Mistral 7B, which applies Proximal Policy Optimization (PPO) training \citep{schulman2017proximal} with a process-level reward model on top of Mistral 7B \citep{mistral}, (7) the WizardMath models \citep{wizardmath}, which strengthen mathematical reasoning in both Mistral 7B and Llama-2 70B \citep{llama2} using evolve-instruct (i.e., instruction tuning with AI-generated instructions) and PPO, primarily leveraging GSM8K and MATH datasets, (8) MetaMath 70B \citep{metamath}, which consists of Llama-2 70B further fine-tuned on expanded GSM8K and MATH data, (9) ToRA 34B \cite{tora}, a CodeLlama 34B variant trained for tool-based mathematical reasoning, and (10) MAmmoTH 70B \citep{MathInstruct}, obtained via instruction tuning Llama-2 70B on MathInstruct.

Analysis presented in Table \ref{tab:sft_rl_math} reveals that when tool assistance is not permitted during evaluation, DeepSeekMath-Instruct 7B achieves high stepwise reasoning performance. Impressively, on the advanced MATH dataset, it outperforms all other open-source models as well as most closed models (e.g., Inflection-2, Gemini Pro) by a margin of at least 9\% in absolute terms. This lead holds even over much larger models such as Qwen 72B or those fine-tuned using math-oriented RL (e.g., WizardMath-v1.1 7B). Although DeepSeekMath-Instruct is on par with Chinese proprietary models like GLM-4 and Baichuan-3 regarding MATH performance, it still trails behind GPT-4 and Gemini Ultra.

When the evaluation setup allows for a combination of natural language reasoning and code/tool invocation, DeepSeekMath-Instruct 7B approaches a 60\% accuracy rate on MATH, outperforming all publicly released models. Across additional benchmarks, it rivals DeepSeek-LLM-Chat 67B—the prior leading model, which possesses tenfold greater scale.

\section{Reinforcement Learning}

\subsection{Group Relative Policy Optimization}
Reinforcement learning (RL) has shown substantial efficacy in refining LLM mathematical reasoning capabilities subsequent to Supervised Fine-Tuning (SFT) \citep{wang2023math,wizardmath}. This section describes our novel and efficient RL approach, termed Group Relative Policy Optimization (GRPO).

\subsubsection{From PPO to GRPO}
Proximal Policy Optimization (PPO) \citep{schulman2017proximal} serves as a widely adopted actor-critic RL algorithm during the RL fine-tuning phase for LLMs \citep{ouyang2022training}. The algorithm tunes LLMs by maximizing the surrogate objective defined as:

\begin{equation}
\footnotesize
    \mathcal{J}_{PPO}(\theta) = \mathbb{E}{[q \sim P(Q), o \sim \pi_{\theta_{old}}(O|q)]} \frac{1}{|o|} \sum_{t=1}^{|o|} \min \left[ \frac{\pi_\theta(o_{t} | q, o_{<t})}{\pi_{\theta_{old}}(o_{t} | q, o_{<t})} A_{t}, \text{clip} \left( \frac{\pi_\theta(o_{t} | q, o_{<t})}{\pi_{\theta_{old}}(o_{t} | q, o_{<t})}, 1 - \varepsilon, 1 + \varepsilon \right)  A_{t} \right]

Here, $\pi_{\theta}$ and $\pi_{\theta_{old}}$ denote the policy models representing the latest and previous iterations, respectively, while $q$ and $o$ are questions and associated outputs, sampled from the dataset of questions and from the historical policy $\pi_{\theta_{old}}$. The parameter $\varepsilon$ serves as the clipping hyperparameter in PPO, providing regularization to ensure stable updates. The term $A_t$ refers to the advantage, which is estimated using Generalized Advantage Estimation (GAE) \citep{gae}, and relies on the reward sequence $\{r_{\ge t}\}$ alongside a parameterized value function $V_{\psi}$. This setup necessitates concurrent training of the value function together with the policy network in PPO. To address excessive reward optimization, it is customary to integrate a per-token KL penalty relative to a reference model into the reward signal at every generation step \citep{ouyang2022training}, which is defined as:

\begin{equation}
    r_{t} = r_\varphi(q, o_{\le t}) - \beta \log\frac{\pi_{\theta}(o_{t}|q, o_{<t})}{\pi_{ref}(o_{t}|q, o_{<t})},
\label{eq:PPO-reward}
\end{equation}

where $r_\varphi$ represents the reward function, $\pi_{ref}$ stands for the reference policy, typically corresponding to the initial SFT model, and $\beta$ denotes the strength of the KL regularization.

Because the value function in PPO is commonly realized as an additional neural model of similar scale to the policy model, this design significantly increases the computational and memory overhead. Furthermore, in reinforcement learning, the value function's role as a baseline is primarily to reduce variance when computing the advantage. For large language models, reward models are usually applied only to the terminal token, assigning it a reward, which poses difficulties for training a value function that is reliable across all generation steps. To overcome these issues, we introduce, as illustrated in Figure \ref{fig:grpo}, Group Relative Policy Optimization (GRPO). GRPO eliminates the requirement for a separate value function by instead adopting the mean reward of a set of outputs—each generated for the same prompt—as the baseline. Specifically, for a given question $q$, GRPO generates a collection of outputs $\{o_1, o_2, \cdots, o_G\}$ from $\pi_{\theta_{old}}$ and optimizes the current policy by maximizing the objective:

\begin{equation}
\footnotesize
\begin{split}
    \mathcal{J}_{GRPO}(\theta) &= \mathbb{E}{[q \sim P(Q), \{o_i\}_{i=1}^G \sim \pi_{\theta_{old}}(O|q)]}  \\
    & \frac{1}{G}\sum_{i=1}^G\frac{1}{|o_i|} \sum_{t=1}^{|o_i|} \left\{ \min \left[ \frac{\pi_\theta(o_{i,t} | q, o_{i,<t})}{\pi_{\theta_{old}}(o_{i,t} | q, o_{i,<t})} \hat{A}_{i,t}, \text{clip} \left( \frac{\pi_\theta(o_{i,t} | q, o_{i,<t})}{\pi_{\theta_{old}}(o_{i,t} | q, o_{i,<t})}, 1 - \varepsilon, 1 + \varepsilon \right)  \hat{A}_{i,t} \right] - \beta \mathbb{D}_{KL}\left[\pi_{\theta} || \pi_{ref}\right]\right\} ,
\end{split}
\label{eq:GRPO-obj}
\end{equation}

In this formulation, $\varepsilon$ and $\beta$ are hyperparameters, and $\hat{A}_{i,t}$ denotes the advantage score that is derived solely based on the relative rewards among outputs within each sampled group, with further explanation provided in subsequent sections. GRPO’s design for advantage computation—using group-wise comparisons—matches the comparative training paradigm common in reward modeling, where reward models are developed on the basis of pairwise preferences for responses to the same question. Moreover, unlike PPO, which injects the KL penalty within the reward term, GRPO applies regularization by explicitly including the KL divergence between the policy and reference distributions in the objective function itself, thereby streamlining the calculation of $\hat{A}_{i,t}$. Furthermore, differing from the KL penalty implemented in (\ref{eq:P

\begin{algorithm}[t]
  \small
  \caption{Iterative Group Relative Policy Optimization}
  \textbf{Input} initial policy model $\pi_{\theta_{\text{init}}}$; reward models $r_\varphi$; task prompts $\mathcal{D}$; 
  hyperparameters $\varepsilon$, $\beta$, $\mu$
  \begin{algorithmic}[1]
    \State Initialize the policy model: $\pi_\theta \leftarrow \pi_{\theta_{\text{init}}}$
    \For{iteration = 1, \dots, I}
      \State Set reference model: $\pi_{ref} \leftarrow \pi_{\theta}$
      \For{step = 1, \dots, M}
      \State Draw a batch $\mathcal{D}_b$ from $\mathcal{D}$
      \State Assign the old policy model: $\pi_{\theta_{old}} \leftarrow \pi_{\theta}$ 
      \State For every question $q$ in $\mathcal{D}_b$, sample $G$ outputs $\{o_i\}_{i=1}^G \sim \pi_{\theta_{old}} (\cdot \mid q) $
      \State Evaluate each output $o_i$ with $r_\varphi$ to compute rewards $\{r_i\}_{i=1}^{G}$
      \State Compute group relative advantages $\hat{A}_{i,t}$ for the $t$-th token of $o_i$
      \For{GRPO iteration = 1, \dots, $\mu$}
        \State Perform policy update on $\pi_{\theta}$ by maximizing the GRPO objective (Equation \ref{eq:GC-GRPO})
      \EndFor
    \EndFor
    \State Update $r_\varphi$ continually by training with a replay strategy. 
    \EndFor 
  \end{algorithmic}
  \textbf{Output} $\pi_\theta$
  \label{alg:iter-grpo}
\end{algorithm}

\subsubsection{Outcome Supervision RL with GRPO} In this setting, for a given question $q$, a set of candidate responses $\{o_1, o_2, \cdots, o_G\}$ is generated using the prior policy model $\pi_{\theta_{old}}$. Each candidate output is then assessed by the reward model to obtain $G$ rewards, $\mathbf{r} = \{r_1, r_2, \cdots, r_G\}$. To standardize the reward values, each reward is centered by subtracting the group mean and scaled by dividing by the group standard deviation. Under outcome supervision, the standardized reward is assigned as the advantage to every token within output $o_i$, specifically, $\hat{A}_{i, t} = \widetilde{r}_i = \frac{r_i- {\rm mean}(\mathbf{r})}{{\rm std}(\mathbf{r})}$. The policy is subsequently optimized by maximizing the objective in equation (\ref{eq:GRPO-obj}).

\subsubsection{Process Supervision RL with GRPO} While outcome supervision only yields a reward at the completion of each output, which may offer limited guidance for complex mathematical reasoning, process supervision is employed to provide feedback after each reasoning step, following the methodology in \cite{wang2023math}. For question $q$ and $G$ candidate outputs $\{o_1, o_2, \cdots, o_G\}$, a process reward model assigns rewards for every step, resulting in $\mathbf{R} = \{ \{r_1^{index(1)},\cdots,r_1^{index(K_1)}\}, \cdots,  \{r_G^{index(1)},\cdots,r_G^{index(K_G)}\} \}$, where $index(j)$ represents the ending token of the $j$-th reasoning step and $K_i$ is the number of steps in output $o_i$. These stepwise rewards are normalized across the group, i.e., $\widetilde{r}_i^{index(j)} = \frac{r_i^{index(j)} - {\rm mean(\mathbf{R})}}{{\rm std(\mathbf{R})}}$. Subsequently, the advantage assigned to each token is determined by summing the normalized rewards from the current token’s position to the end, that is, $\hat{A}_{i, t} = \sum_{index(j) \ge t} \widetilde{r}_i^{index(j)}$. The policy parameters are then updated to maximize the objective as specified in equation (\ref{eq:GRPO-obj}).

\subsubsection{Iterative RL with GRPO} As training proceeds in reinforcement learning, reliance on an earlier reward model may become inadequate for supervising the latest policy model. To address this, we examine an iterative RL scheme leveraging GRPO. In this setup (illustrated in Algorithm \ref{alg:iter-grpo}), we construct updated reward model training datasets by sampling new data from the evolving policy model. The reward model is updated continuously via a replay approach, maintaining 10\% of its previous training data to preserve historical context. Next, the policy model is aligned with the reward model by resetting the reference model to the policy model itself and refining the policy with the improved reward signal.

\subsection{Training and Evaluating DeepSeekMath-RL}

For reinforcement learning, we utilize DeepSeekMath-Instruct 7B as our base. The RL training set comprises approximately 144K chain-of-thought-formatted questions sourced from GSM8K and MATH segments within the SFT dataset; other SFT-derived questions are deliberately excluded to isolate RL’s effects on benchmarks lacking training data during the RL procedure. Reward model training data construction follows the methodology described in \citep{wang2023math}. Our initial reward model is trained using DeepSeekMath-Base 7B, adopting a learning rate of 2e-5. In GRPO, we apply a policy model learning rate of 1e-6 and set the KL regularization coefficient at 0.04. For each item, $64$ samples are produced. The maximum sequence length is 1024, with a batch size of 1024. After each exploratory sampling phase, the policy model is updated once. Evaluation of DeepSeekMath-RL 7B follows the same protocol as DeepSeekMath-Instruct 7B, where GSM8K and MATH (with chain-of-thought reasoning) are considered in-domain tasks, and all other benchmarks serve as out-of-domain evaluations.

Table \ref{tab:sft_rl_math} presents a comparative analysis of the performance of open- and closed-source models, evaluating both chain-of-thought and tool-augmented reasoning across English and Chinese datasets. The key observations are as follows: 1) DeepSeekMath-RL 7B achieves accuracy rates of 88.2\% on GSM8K and 51.7\% on MATH by leveraging chain-of-thought reasoning. This result not only leads the open-source models within the 7B–70B parameter spectrum but also exceeds most closed-source alternatives. 2) Notably, DeepSeekMath-RL 7B undergoes instruction tuning exclusively on chain-of-thought-formatted data from GSM8K and MATH, commencing from the DeepSeekMath-Instruct 7B checkpoint. Despite this relatively narrow training dataset, it consistently outperforms DeepSeekMath-Instruct 7B across all benchmarks, underscoring the benefits conferred by reinforcement learning.

\section{Discussion}
In this section, we detail insights gained from our work on both pre-training and reinforcement learning phases.

\subsection{Lessons Learnt in Pre-Training}

We begin by recounting key takeaways from the pre-training process. Unless specified otherwise, the training configurations follow those outlined in Section \ref{sec:quality-policy}. It is important to clarify that the term DeepSeekMath Corpus, as used here, refers specifically to the 89B-token dataset curated during the second round of data collection.

\subsubsection{Code Training Benefits Mathematical Reasoning}

A widely held, yet not fully substantiated, conjecture posits that pre-training on code improves a model's reasoning capabilities. Our work contributes partial evidence in support of this, particularly with respect to mathematical problem solving: training on code substantially enhances models' proficiency in mathematical reasoning, with and without access to computational tools.

To assess the impact of code-centric training on mathematical reasoning performance, we evaluated both a two-stage and a one-stage training regimen:

\textbf{Two-Stage Training}

\begin{itemize}[topsep=0pt]
    \item \textbf{Code Training for 400B Tokens $\rightarrow$ Math Training for 150B Tokens}: DeepSeek-LLM 1.3B is first pretrained on 400B code tokens, after which it undergoes an additional 150B tokens of math-specific training;
    \item \textbf{General Training for 400B Tokens $\rightarrow$ Math Training for 150B Tokens}: To provide a baseline for comparison, we also replace the code token pretraining phase with 400B general domain tokens (sourced from DeepSeek-AI’s extensive general corpus), and follow with 150B math tokens. This allows us to discern the relative contribution of code versus general pretraining on mathematical reasoning capabilities.
\end{itemize}

\textbf{One-Stage Training}

\begin{itemize}[topsep=0pt]
    \item \textbf{Math Training for 150B Tokens}: The model is directly trained on 150B tokens consisting exclusively of mathematical content;
    \item \textbf{Training on a mixture of 400B Code Tokens and 150B Math Tokens}: Since sequential code and math training can impair coding skills (catastrophic forgetting), we assess whether exposing the model to a combined dataset—integrating code and math tokens during a unified training phase—can concurrently strengthen mathematical reasoning and maintain programming proficiency.
\end{itemize}

\paragraph{Results}
The downstream task results for these various training approaches are summarized in Table \ref{tab:code-math} and Table \ref{tab:code-math-general-eval}.

Introducing code training proves advantageous for mathematical reasoning augmented by program execution, regardless of whether the model undergoes two-stage or one-stage training. Evidence from Table \ref{tab:code-math} indicates that models pretrained on code substantially outperform others in Python-based mathematical problem solving (e.g., GSM8K and MATH datasets) in two-stage regimens. An extra phase of math training provides further enhancement. Notably, in the one-stage paradigm, simultaneously mixing code and math data reduces catastrophic forgetting observed with sequential approaches and produces gains for both coding ability (see Table \ref{tab:code-math-general-eval}) and math reasoning that leverages programming (see Table \ref{tab:code-math}).

Code-oriented pretraining also supports mathematical reasoning even when no external tools are available. In the two-stage scenario, beginning with code yields moderate improvements and accelerates learning in subsequent math training, resulting in the top observed scores. In contrast, training with an interleaved code/math mixture in a single stage appears to weaken non-tool-based mathematical reasoning. This may be attributed to the limited capacity of DeepSeek-LLM 1.3B, which potentially cannot effectively learn from the combined data distributions when trained in parallel.

\subsubsection{ArXiv Papers Appear Limited in Enhancing Mathematical Reasoning}

It is standard practice to include arXiv papers as part of mathematical pre-training corpora \citep{minerva,gpt-f,llemma,mathpile}. Nonetheless, their actual contribution to the development of mathematical reasoning capabilities has not been thoroughly scrutinized. Unexpectedly, our findings indicate that incorporating arXiv papers does not significantly enhance mathematical reasoning. To examine this, we trained models across different scales, specifically DeepSeek-LLM 1.3B and DeepSeek-Coder-Base-v1.5 7B \citep{deepseek-coder}, using variously processed arXiv datasets:

\begin{itemize}[topsep=0pt]
    \item \textbf{MathPile} \citep{mathpile}: a dataset comprising 8.9B tokens processed using cleaning and heuristic filtering strategies, with more than 85\% of the content sourced from scientific arXiv papers;
    \item \textbf{ArXiv-RedPajama} \citep{redpajama}: the full set of arXiv LaTeX documents, with preambles, comments, macros, and bibliographies eliminated, amounting to 28.0B tokens.
\end{itemize}

Our experiments involved training DeepSeek-LLM 1.3B for 150B tokens and DeepSeek-Coder-Base-v1.5 7B for 40B tokens on each respective arXiv dataset. The outcome suggests that using arXiv-based corpora yields little to no discernible advantage in mathematical reasoning performance. In fact, when evaluated on a variety of mathematical benchmarks at diverse difficulty levels, both models trained solely on arXiv material did not exhibit meaningful gains, and sometimes even underperformed. The assessed benchmarks include quantitative reasoning tasks like GSM8K and MATH (refer to Table \ref{tab:arxiv-cot}), multiple-choice assessments such as MMLU-STEM (see Table \ref{tab:arxiv-cot}), as well as formal mathematics evaluation like miniF2F (see Table \ref{tab:arxiv_atp}).

That said, these observations are subject to certain caveats. Specifically, the following remain unexamined in our study:

\begin{itemize}[topsep=0pt]
    \item The effects of arXiv-based training on specialized mathematics tasks not explored here, such as informalizing theorems—i.e., transforming formal proofs or statements into more informal explanations;
    \item The potential influence of arXiv-derived data in combination with alternative datasets;
    \item Whether scaling up model parameters could eventually reveal advantages attributable to arXiv paper pre-training.
\end{itemize}

Therefore, further comprehensive investigations are warranted, which are left open for subsequent research efforts.

\subsection{Insights of Reinforcement Learning}
\subsubsection{Toward a Unified Paradigm} In what follows, we propose a framework that brings together a range of training strategies, including SFT, RFT, DPO, PPO, GRPO. We also empirically examine how different elements of this unified paradigm influence training. More broadly, the gradient of a training objective with respect to parameter $\theta$ can be formalized as follows:

\begin{equation}
    \nabla_{\theta}\mathcal{J}_{\textcolor{red}{\mathcal{A}}}(\theta) = \mathbb{E}[\underbrace{(q,o) \sim \textcolor{red}{\mathcal{D}}}_{Data \ Source}]\left( \frac{1}{|o|} \sum_{t=1}^{|o|}  \underbrace{GC_{{\mathcal{A}}}(q, o, t, \textcolor{red}{\pi_{{rf}}})}_{Gradient \ Coefficient}  \nabla_{\theta}\log \pi_{\theta}(o_t | q, o_{<t})\right).
\label{eq:objective}
\end{equation}

The equation above relies on three principal elements: (1) \textit{Data Source $\mathcal{D}$}, which specifies the set of training instances; (2) \textit{Reward Function $\pi_{{rf}}$}, which provides the feedback signal for learning; and (3) \textit{Algorithm $\mathcal{A}$}, responsible for transforming both training instances and reward signals into the gradient coefficient $GC$, ultimately determining the scale of reinforcement or penalization applied. We interpret several key training strategies within this shared framework:

\begin{itemize}
    \item  \textbf{Supervised Fine-tuning (SFT)}: In SFT, the pretrained model undergoes further training using curated, human-labeled data.
    \item \textbf{Rejection Sampling Fine-tuning (RFT)}: This method fine-tunes an SFT-trained model with outputs from the same model, filtering these outputs based on answer correctness for a given SFT prompt.
    \item \textbf{Direct Preference Optimization (DPO)}: DPO improves upon the SFT model by incorporating additional, sampled outputs and optimizing with a pairwise DPO loss.
    \item \textbf{Online Rejection Sampling Fine-tuning (Online RFT)}: Distinct from RFT, Online RFT begins with an SFT-initialized policy and iteratively refines it using augmented samples generated from the current, up-to-date policy.
    \item \textbf{PPO/GRPO}: These methods initialize the policy model from the SFT baseline, then enhance performance by reinforcement learning on outputs dynamically sampled from the live policy model.
\end{itemize}

A comprehensive overview of these methods and their components is provided in Table \ref{tab:data_policy}. For an in-depth explanation of the derivations involved, see Appendix \ref{app:analysis-rl}.

\paragraph{Analysis of Data Source} We classify data sources into two main types: online sampling and offline sampling. Online sampling involves collecting training samples generated by the currently evolving policy model through exploration, whereas offline sampling utilizes data obtained from the output of the original SFT model. Both RFT and DPO employ the offline sampling paradigm, while Online RFT and GRPO make use of online sampling.

As illustrated in Figure \ref{fig:rl-analysis}, our findings reveal that Online RFT notably surpasses RFT in performance across both benchmark tasks. At the early phase of training, Online RFT achieves results similar to those of RFT. However, as training progresses, Online RFT gains a distinct advantage, underscoring the benefits of an online training framework. This outcome is expected, as the policy model and the SFT model are initially similar, resulting in minimal divergence within the sampled data. Over the course of training, however, data sampled from the evolving policy (the actor) becomes increasingly distinct, with online sampling conferring substantial performance improvements.

\paragraph{Analysis of Gradient Coefficient} The algorithms leverage reward signals to shape the gradient coefficient, thereby guiding parameter updates. In our empirical setup, the reward function is categorized as either 'Rule' or 'Model.' 'Rule' corresponds to assessing responses based on explicit correctness, whereas 'Model' involves utilizing a reward model—trained using rule-based evaluations—to assign scores to responses. Equations \ref{eq:GC-RFT} and \ref{eq:GC-GRPO} elucidate a fundamental distinction: GRPO dynamically adjusts the gradient coefficient according to reward magnitudes from the reward model, enabling differentiated reinforcement and penalization depending on the response quality. By contrast, Online RFT does not possess such adaptability; it refrains from penalizing incorrect answers and treats all correct answers with uniform reinforcement intensity.

As illustrated in Figure \ref{fig:rl-analysis}, GRPO outperforms online RFT, underscoring the advantage of modulating the coefficients for positive and negative gradients. Additionally, GRPO+PS achieves better results than GRPO+OS, which suggests that employing more granular, step-sensitive gradient coefficients is beneficial. Beyond this, we investigate the effects of iterative reinforcement learning by executing two consecutive iterations in our experimental framework. Results, presented in Figure \ref{fig:iter-rl}, reveal that the adoption of iterative RL leads to noticeable gains, particularly after the initial iteration.

\subsubsection{Why RL Works?} In the present work, we apply reinforcement learning to a subset of the instruction tuning dataset, yielding substantial improvements over the base instruction-tuned model. To elucidate the mechanism behind this enhancement, we assess Pass@K and Maj@K accuracy for both Instruct and RL-optimized models across two standard benchmarks. The results, visualized in Figure \ref{fig:maj-pass}, indicate that while RL drives improvements in Maj@K, it does not boost Pass@K. This suggests that reinforcement learning augments the model’s robustness by making the output distribution more reliable—specifically, \textbf{the gains stem from raising the correct response within the TopK outputs rather than fundamentally strengthening baseline capabilities.} An analogous issue, namely the \textbf{misalignment problem} in SFT models for reasoning tasks, was observed in \citep{wang2023making}, where performance in reasoning was enhanced through a suite of preference alignment strategies \citep{yuan2023rrhf,song2023preference,wang2023making}.

\subsubsection{How to Achieve More Effective RL?}
\label{sec:effec_rl}
Our experiments confirm that RL can be highly effective for mathematical reasoning challenges. We further offer a generalized framework to interpret multiple representative training approaches, treating them as variants or simplifications of the broader RL methodology. According to Equation \ref{eq:objective}, three principal elements characterize this framework: Data Source, Algorithm, and Reward Function. We outline prospective research avenues pertaining to each component.

\paragraph{Data Source} The data source is fundamental to every training method. For RL, this entails unlabeled questions and their corresponding outputs sampled from the policy model. This study employs question data from the instruction tuning phase, alongside straightforward nucleus sampling to generate outputs. We hypothesize that this limitation partly explains why our RL process boosts Maj@K exclusively. In future research, we plan to extend our RL approach to out-of-distribution prompts and incorporate \textbf{advanced sampling (decoding) strategies}, such as those inspired by tree-search methods \citep{yao2023tree}. Moreover, the deployment of \textbf{efficient inference techniques} \citep{xia-etal-2023-speculative,leviathan2023fast,kwon2023efficient,xia2024unlocking}—which shape how effectively policy models explore—will also be crucial for improving the overall efficiency of exploration.

\paragraph{Algorithms} Algorithms utilize both data and reward signals to compute gradient coefficients, which in turn adjust the model parameters. According to Equation \ref{eq:objective}, most current approaches implicitly place complete \textbf{TRUST} in the feedback provided by the reward function, using it to modulate the conditional probability of selecting specific tokens. Nonetheless, the reliability of the reward signal cannot always be guaranteed, particularly in the context of highly intricate tasks. A case in point is the PRM800K dataset \citep{lightman2023let}, which, despite being thoroughly annotated by skilled experts, is reported to have around 20\% erroneous labels\footnote{\url{https://github.com/openai/prm800k/issues/12\#issuecomment-1728491852}}. Addressing these challenges, our focus is directed toward reinforcement learning algorithms capable of withstanding noisy or imperfect reward feedback. We posit that alignment approaches such as \textbf{WEAK-TO-STRONG} \citep{burns2023weak} are likely to catalyze profound advancements in learning algorithm design.

\paragraph{Reward Function} The reward function serves as the foundation for the learning signal in RL, typically instantiated as a neural reward model. We identify three pivotal research avenues for reward models: 1) \textbf{Strategies for improving reward model generalization.} It is crucial for reward models to transfer effectively to unseen queries and sophisticated generated outputs; without this, RL may only consolidate the present output distribution of LLMs rather than significantly elevate their underlying performance; 2) \textbf{Capturing and utilizing reward model uncertainty.} Quantifying uncertainty may enable new connections between less robust reward functions and weak-to-strong learning paradigms; 3) \textbf{Developing efficient and high-quality process reward models} capable of delivering detailed supervision signals during step-by-step reasoning procedures \citep{lightman2023let,wang2023math}.

\section{Conclusion, Limitation, and Future Work} In this work, we introduced DeepSeekMath, which establishes state-of-the-art results among open-source models on the competition-grade MATH benchmark and achieves performance levels comparable to proprietary models. DeepSeekMath~builds on DeepSeek-Coder-v1.5 7B and is further continually trained over 500B tokens, with 120B math-focused tokens curated from Common Crawl comprising a substantial portion of the dataset. Our thorough ablation analyses reveal that online web pages constitute a valuable resource for acquiring quality mathematical data, whereas content from arXiv does not contribute as significantly as initially anticipated. We present Group Relative Policy Optimization (GRPO), an adaptation of Proximal Policy Optimization (PPO), that can substantially bolster mathematical reasoning proficiency while optimizing for lower memory requirements. Experimental results demonstrate that GRPO continues to provide substantial benefits even when DeepSeekMath-Instruct 7B already performs at a strong level on various benchmarks. Additionally, we offer a cohesive framework that unifies a family of methods and outline several prospective avenues for advancing reinforcement learning techniques.

Despite these advances, DeepSeekMath~exhibits comparatively weaker performance in tasks related to geometry and formal theorem proving compared to closed-source counterparts. Specifically, in preliminary evaluations, the model struggles with certain types of geometry questions—such as those involving triangles and ellipses—potentially due to limitations or imbalances in the data selection strategy used during pre-training and fine-tuning. Furthermore, given the model’s parameter size, DeepSeekMath~lags behind GPT-4 in few-shot learning scenarios; GPT-4 shows notable improvement with few-shot examples, while DeepSeekMath~yields nearly equivalent outcomes in both zero-shot and few-shot tests. Looking ahead, we aim to further refine our engineered data selection processes to assemble an even higher quality pre-training corpus. We will also investigate the research directions described in Section \ref{sec:effec_rl} to enable more effective reinforcement learning for LLMs.

\end{document}